\subsection{有序群的孤立子群}

为了从赋值的观点研究第 1 节中的情形，我们需要下面的定义 1 和命题 3。

定义 1. 有序群 G 的子群 H 称为孤立的，如果关系 \(0 \leqslant y \leqslant x\) 与 \(x \in \mathbf{H}\) 蕴含 \(y \in \mathbf{H}\) 。

例 (1) 设 A 与 B 是两个有序群；给 \(A \times B\) 赋以字典序（即~“ \((a, b) \leqslant (a', b')\) ”等价于“ \((a < a')\) 或 \((a = a' \text{ 且 } b \leqslant b')\) ”）。立即可见，\(A \times B\) 的第二个因子 B 是 \(A \times B\) 的孤立子群。

命题 3. 设 G 是有序群，P 是其正元素集合。

\begin{enumerate}
\def\labelenumi{(\alph{enumi})}
\item
  由 G 到某个有序群的递增同态的核是 G 的孤立子群。
\item
  反之，设 H 是 G 的孤立子群，g 是 G 到 G/H 上的典范同态。那么 \(g(P)\) 是 G/H 上某个有序群结构的正元素集合。此外，若 G 是全序的，则 G/H 也是全序的。
\item
  设 \(f\) 是由 \(G\) 到某个有序群的递增同态；设 \(H\) 表示 \(f\) 的核。若 \(0 \leqslant y \leqslant x\) 且 \(x \in H\) ，则 \(0 \leqslant f(y) \leqslant f(x) = 0\) ，于是 \(f(y) = 0\) ，即 \(y \in H\) 。因此 \(H\) 是孤立的。
\item
  设 \(\mathbf{H}\) 是 \(\mathbf{G}\) 的孤立子群，\(g\colon \mathbf{G}\to \mathbf{G} / \mathbf{H}\) 。设 \(\mathbf{P}' = g(\mathbf{P})\) 。显然 \(\mathbf{P}' + \mathbf{P}'\subset \mathbf{P}'\) 。并且
\end{enumerate}

\[
\mathrm{P} ^ {\prime} \cap (- \mathrm{P} ^ {\prime}) = \{0 \},
\]

因为，若 \(x\) 与 \(y\) 是 \(\mathbf{P}\) 中满足 \(g(x) = -g(y)\) 的元素，则 \(x + y \in \mathbf{H}\) ，于是 \(x \in \mathbf{H}\) 且 \(y \in \mathbf{H}\) ，因为 \(\mathbf{H}\) 是孤立的；因此 \(g(x) = g(y) = 0\) 。这样，\(\mathbf{P}'\) 是 \(\mathbf{G} / \mathbf{H}\) 上某个有序群结构的正元素集合（《代数》第六章 §1 第 3 节命题 3）。最后，若 \(\mathbf{G}\) 是全序的，则 \(\mathbf{P} \cap (-\mathbf{P}) = \mathbf{G}\) ，于是 \(\mathbf{P}' \cup (-\mathbf{P}') = \mathbf{G} / \mathbf{H}\) ，因此 \(\mathbf{G} / \mathbf{H}\) 是全序的（同前）。

例 (2) 若重新考察 G 是字典积 \(A \times B\) 且 H = B 的例子，则有序群 G/H 典范地等同于 A。

